%============================================================
% CH1 · LESSON 1 — VERTICES AND EDGES
%============================================================
\subsection{Vertices and Edges}
\label{ssec:ch1-vertices-edges}

\subsubsection*{The one-paragraph answer}

A graph is the simplest possible description of \emph{things and the connections between them}. The things are called \textbf{vertices} (or nodes), the connections are called \textbf{edges}. Cities and roads, people and friendships, web pages and links, courses and prerequisites, cells of a maze and the doors between them --- all of these are graphs. The power of the idea is that once you have drawn the dots and the lines, the \emph{meaning} of the dots no longer matters: the same few algorithms answer ``can I get from here to there?'', ``how many steps does it take?'', ``which groups fall apart?'' for roads, friendships and mazes alike.

\begin{definitionbox}[Definition (Graph)]
A \textbf{graph} $G=(V,E)$ consists of
\begin{itemize}[leftmargin=*]
    \item a finite set $V$ of \textbf{vertices}; we write $n=\nV$ for their number;
    \item a set $E$ of \textbf{edges}; each edge connects two vertices $u$ and $v$, its \textbf{endpoints}. We write $m=\nE$.
\end{itemize}
Two vertices joined by an edge are \textbf{adjacent} (or \textbf{neighbours}); the edge is \textbf{incident} to both of them. In contest input the vertices are almost always the integers $1,2,\dots,n$ and the edges arrive as a list of pairs.
\end{definitionbox}

\subsubsection*{A first picture: five towns}

Five towns A, B, C, D, E and the roads between them. The geography is gone; only ``who is connected to whom'' remains.

\begin{diagrambox}
\begin{tikzpicture}[x=1.3cm, y=1.1cm]
    % left: the "map"
    \begin{scope}
        \fill[insightbg, rounded corners=6pt] (-0.6,-1.7) rectangle (3.6,1.5);
        \node[font=\scriptsize\bfseries\color{darkgray}] at (1.5,1.25) {a map};
        \foreach \p/\n in {(0,0.6)/A, (1.3,0.9)/B, (2.9,0.5)/C, (0.6,-1.1)/D, (2.4,-1.0)/E}
            \node[rectangle, draw=darkgray, fill=white, font=\scriptsize\bfseries, inner sep=2pt] (m\n) at \p {\n};
        \draw[accentorange, line width=1.6pt, rounded corners=4pt] (mA) -- (0.7,0.9) -- (mB);
        \draw[accentorange, line width=1.6pt] (mB) .. controls (2.2,1.1) .. (mC);
        \draw[accentorange, line width=1.6pt] (mA) .. controls (0.1,-0.5) .. (mD);
        \draw[accentorange, line width=1.6pt] (mD) -- (mE);
        \draw[accentorange, line width=1.6pt] (mB) .. controls (1.6,-0.2) .. (mE);
        \draw[accentorange, line width=1.6pt] (mC) -- (mE);
    \end{scope}
    \draw[flow, line width=1.2pt] (3.9,0) -- node[above, font=\scriptsize] {keep only} node[below, font=\scriptsize] {connections} (5.3,0);
    % right: the graph
    \begin{scope}[xshift=7.2cm]
        \node[gv] (A) at (0,0.6) {A};
        \node[gv] (B) at (1.3,1.0) {B};
        \node[gv] (C) at (2.6,0.5) {C};
        \node[gv] (D) at (0.4,-1.1) {D};
        \node[gv] (E) at (2.0,-1.0) {E};
        \foreach \a/\b in {A/B, B/C, A/D, D/E, B/E, C/E} \draw[ge] (\a) -- (\b);
        \node[font=\scriptsize\bfseries\color{darkgray}] at (1.3,1.6) {a graph: $n=5$, $m=6$};
    \end{scope}
\end{tikzpicture}
\figcap{The road map becomes a graph with $V=\{A,B,C,D,E\}$ and $E=\{AB, BC, AD, DE, BE, CE\}$. Lengths, curves and positions are forgotten: the picture on the right could be redrawn in any layout and it would still be the same graph.}
\end{diagrambox}

\begin{insightbox}[The drawing is not the graph]
A graph is only the two sets $V$ and $E$. Where you put the dots and whether the lines cross is your choice; two very different-looking pictures can show the same graph. Draw the version that makes the structure obvious --- this is a skill, and every chapter of this book practises it.
\end{insightbox}

%------------------------------------------------------------
\subsubsection*{Undirected and directed edges}

A road usable in both directions is an \textbf{undirected} edge $\{u,v\}$: it can be walked from $u$ to $v$ and from $v$ to $u$. A one-way street, a ``follows'' relation on a social network or ``course $u$ must be taken before course $v$'' is a \textbf{directed} edge (an \textbf{arc}) $u\to v$: an ordered pair $(u,v)$ with a \textbf{tail} $u$ and a \textbf{head} $v$. A graph whose edges are all directed is a \textbf{directed graph} (\textbf{digraph}).

\begin{diagrambox}
\begin{tikzpicture}[x=1.25cm, y=1.1cm]
    \begin{scope}
        \node[gv] (1) at (0,1) {1};
        \node[gv] (2) at (1.6,1) {2};
        \node[gv] (3) at (0,-0.3) {3};
        \node[gv] (4) at (1.6,-0.3) {4};
        \foreach \a/\b in {1/2, 1/3, 2/3, 2/4, 3/4} \draw[ge] (\a) -- (\b);
        \node[font=\small\bfseries] at (0.8,-1.15) {undirected};
        \node[font=\scriptsize\color{darkgray}] at (0.8,-1.6) {friendship, two-way road};
    \end{scope}
    \begin{scope}[xshift=4.2cm]
        \node[gv] (1) at (0,1) {1};
        \node[gv] (2) at (1.6,1) {2};
        \node[gv] (3) at (0,-0.3) {3};
        \node[gv] (4) at (1.6,-0.3) {4};
        \foreach \a/\b in {1/2, 3/1, 2/3, 2/4, 4/3} \draw[gde] (\a) -- (\b);
        \node[font=\small\bfseries] at (0.8,-1.15) {directed};
        \node[font=\scriptsize\color{darkgray}] at (0.8,-1.6) {follows, one-way street};
    \end{scope}
    \begin{scope}[xshift=8.4cm]
        \node[gv] (1) at (0,1) {1};
        \node[gv] (2) at (1.6,1) {2};
        \draw[ge] (1) -- (2);
        \node[font=\scriptsize] at (0.8,0.35) {is stored as};
        \node[gv] (1b) at (0,-0.3) {1};
        \node[gv] (2b) at (1.6,-0.3) {2};
        \draw[gde] (1b) to[bend left=25] (2b);
        \draw[gde] (2b) to[bend left=25] (1b);
        \node[font=\small\bfseries] at (0.8,-1.15) {in code};
        \node[font=\scriptsize\color{darkgray}] at (0.8,-1.6) {one undirected = two arcs};
    \end{scope}
\end{tikzpicture}
\figcap{Left: from 2 you can go to 1, 3 or 4. Middle: from 2 you can go to 3 and 4 but \emph{not} to 1 (the arc is $1\to 2$). Right: programs store an undirected edge as two opposite arcs (Lesson~\ref{ssec:ch1-storing}).}
\end{diagrambox}

\begin{warningbox}[Read the statement for the direction]
``There is a road between $u$ and $v$'' is undirected; ``$u$ trusts $v$'', ``$u$ points to $v$'', ``from room $u$ you can enter room $v$'', ``$a_i$ is the next element after $i$'' are directed. Treating a directed edge as undirected (or forgetting the reverse copy of an undirected edge) is one of the most frequent wrong answers in graph problems.
\end{warningbox}

%------------------------------------------------------------
\subsubsection*{Weighted edges}

Often an edge carries a number: a road length, a ticket price, a travel time, a capacity. Such a graph is \textbf{weighted}, and we write $w(u,v)$ for the weight of edge $uv$. An \textbf{unweighted} graph behaves as if every edge had weight 1, so ``shortest'' then means ``fewest edges''.

\begin{diagrambox}
\begin{tikzpicture}[x=1.5cm, y=1.1cm]
    \node[gvsrc] (s) at (0,0) {S};
    \node[gv] (a) at (1.4,1) {A};
    \node[gv] (b) at (1.4,-1) {B};
    \node[gvdst] (t) at (2.8,0) {T};
    \draw[ge] (s) -- node[gw, above left] {2} (a);
    \draw[ge] (s) -- node[gw, below left] {5} (b);
    \draw[ge] (a) -- node[gw, above right] {7} (t);
    \draw[ge] (b) -- node[gw, below right] {1} (t);
    \draw[ge] (a) -- node[gw, right] {1} (b);
    \begin{scope}[xshift=5.5cm]
        \node[gvsrc] (s) at (0,0) {S};
        \node[gv] (a) at (1.4,1) {A};
        \node[gv] (b) at (1.4,-1) {B};
        \node[gvdst] (t) at (2.8,0) {T};
        \draw[gpath] (s) -- (a); \draw[gpath] (a) -- (b); \draw[gpath] (b) -- (t);
        \draw[gfaded] (s) -- (b); \draw[gfaded] (a) -- (t);
        \node[gw] at (0.6,0.65) {2}; \node[gw] at (1.62,0) {1}; \node[gw] at (2.2,-0.65) {1};
        \node[gdist, right=1mm of t] {$2+1+1=4$};
    \end{scope}
\end{tikzpicture}
\figcap{A weighted graph. The cheapest route from \legendchip{gvsrc}{S} to \legendchip{gvdst}{T} is S--A--B--T of total weight 4, although it uses three edges; the routes S--B--T (weight 6) and S--A--T (weight 9) use only two. Weights change the answer --- shortest paths with weights are the subject of Chapter~7.}
\end{diagrambox}

%------------------------------------------------------------
\subsubsection*{Simple graphs, multigraphs and self-loops}

Two more kinds of edges can appear in real input:
\begin{itemize}[leftmargin=*]
    \item a \textbf{self-loop} is an edge from a vertex to itself, $u$--$u$ (``room 3 has a door back into room 3'');
    \item \textbf{parallel edges} (multi-edges) are two or more edges between the same pair of vertices (two different roads between the same towns, maybe of different lengths).
\end{itemize}
A graph with neither is \textbf{simple}; a graph that allows them is a \textbf{multigraph}.

\begin{diagrambox}
\begin{tikzpicture}[x=1.25cm, y=1.1cm]
    \begin{scope}
        \node[gv] (1) at (0,1) {1};
        \node[gv] (2) at (1.6,1) {2};
        \node[gv] (3) at (1.6,-0.3) {3};
        \node[gv] (4) at (0,-0.3) {4};
        \foreach \a/\b in {1/2, 2/3, 3/4, 4/1} \draw[ge] (\a) -- (\b);
        \node[font=\small\bfseries] at (0.8,-1.15) {simple};
    \end{scope}
    \begin{scope}[xshift=5cm]
        \node[gv] (1) at (0,1) {1};
        \node[gv] (2) at (1.6,1) {2};
        \node[gvcur] (3) at (1.6,-0.3) {3};
        \node[gv] (4) at (0,-0.3) {4};
        \draw[ge] (1) -- (2);
        \draw[ge, draw=recursionred, line width=1.3pt] (2) to[bend left=22] (3);
        \draw[ge, draw=recursionred, line width=1.3pt] (2) to[bend right=22] (3);
        \draw[ge] (3) -- (4); \draw[ge] (4) -- (1);
        \draw[ge, draw=recursionred, line width=1.3pt] (3) to[out=-20, in=-80, looseness=7] (3);
        \node[font=\scriptsize\color{recursionred}, anchor=west] at (2.0,1.0) {parallel edges 2--3};
        \node[font=\scriptsize\color{recursionred}, anchor=west] at (2.3,-0.85) {self-loop 3--3};
        \node[font=\small\bfseries] at (0.8,-1.15) {multigraph};
    \end{scope}
\end{tikzpicture}
\figcap{Right: the input \texttt{1 2 / 2 3 / 3 2 / 3 3 / 3 4 / 4 1} ($n=4$, $m=6$). The edges \texttt{2 3} and \texttt{3 2} are the same undirected pair, so they are parallel; \texttt{3 3} is a self-loop.}
\end{diagrambox}

Most contest statements promise ``no self-loops and no multiple edges''. When they do \emph{not}, the promise is missing for a reason: the tests will contain them. Most algorithms in this book work on multigraphs unchanged, but counting problems (``how many neighbours does $v$ have?'') and adjacency matrices (one cell per pair) do not. The following program detects both kinds; it is a good first exercise in reading a graph.

\begin{lstlisting}[style=cpp]
#include <bits/stdc++.h>
using namespace std;

// Reads an undirected graph and reports whether it is simple.
// Input: n m, then m lines "u v" (1-indexed).
int main() {
    int n, m;
    if (!(cin >> n >> m)) return 0;
    set<pair<int,int>> seen;          // each edge stored as (smaller, larger)
    int loops = 0, parallel = 0;
    for (int i = 0; i < m; i++) {
        int u, v;
        cin >> u >> v;
        if (u == v) { loops++; continue; }        // self-loop u-u
        if (u > v) swap(u, v);                    // {u,v} == {v,u}
        if (!seen.insert(make_pair(u, v)).second) // already there?
            parallel++;                           // a second copy
    }
    cout << "self-loops: " << loops << "\n";
    cout << "parallel edges: " << parallel << "\n";
    cout << (loops == 0 && parallel == 0 ? "simple graph" : "multigraph") << "\n";
    return 0;
}
\end{lstlisting}

\begin{lstlisting}[style=trace]
input : 4 6   1 2   2 3   3 2   3 3   3 4   4 1
output: self-loops: 1
        parallel edges: 1
        multigraph
\end{lstlisting}

The line \texttt{if (u > v) swap(u, v)} deserves a name: \textbf{normalising} an undirected edge so that $\{u,v\}$ and $\{v,u\}$ look identical. You will use it whenever undirected edges go into a \texttt{set} or a \texttt{map}.

%------------------------------------------------------------
\subsubsection*{Graphs in real life and in contests}

\begin{table}[H]
\centering\small
\begin{tabularx}{\linewidth}{@{}l l l X@{}}
\toprule
\textbf{Situation} & \textbf{Vertices} & \textbf{Edges} & \textbf{Kind and typical question} \\
\midrule
Road network & towns, crossings & roads & undirected, weighted: shortest route (Ch.~7) \\
Social network & people & friendships / follows & undirected / directed: groups (Ch.~4) \\
The web & pages & hyperlinks & directed: which pages are reachable \\
Course plan & courses & ``$u$ before $v$'' & directed, no cycles: a valid order (Ch.~5) \\
Maze, grid map & free cells & shared sides & undirected, implicit: flood fill (Pattern~3) \\
Puzzle, lock & configurations & one move & implicit: fewest moves (Pattern~4) \\
Permutation $p$ & indices $i$ & $i\to p_i$ & one outgoing arc per vertex: cycles (Pattern~5) \\
Company & employees & boss $\to$ subordinate & a rooted tree (Ch.~13) \\
\bottomrule
\end{tabularx}
\end{table}

\begin{diagrambox}
\begin{tikzpicture}[x=1cm, y=1cm]
    % social
    \begin{scope}
        \node[gvA, minimum size=8.5mm] (a) at (0,0.8) {\scriptsize Ann};
        \node[gvA, minimum size=8.5mm] (b) at (1.4,1.2) {\scriptsize Bob};
        \node[gvA, minimum size=8.5mm] (c) at (1.7,-0.1) {\scriptsize Cy};
        \node[gvB, minimum size=8.5mm] (d) at (0.1,-0.7) {\scriptsize Dee};
        \node[gvB, minimum size=8.5mm] (e) at (1.5,-1.4) {\scriptsize Eve};
        \draw[ge] (a) -- (b); \draw[ge] (b) -- (c); \draw[ge] (a) -- (c);
        \draw[ge] (d) -- (e);
        \node[font=\scriptsize\bfseries] at (0.8,-2.2) {friends: two groups};
    \end{scope}
    % prerequisites
    \begin{scope}[xshift=4.2cm]
        \node[gv, minimum size=9mm] (p) at (0,1) {\scriptsize Prog};
        \node[gv, minimum size=9mm] (d) at (1.9,1) {\scriptsize DS};
        \node[gv, minimum size=9mm] (m) at (0,-0.8) {\scriptsize Math};
        \node[gv, minimum size=9mm] (a) at (1.9,-0.8) {\scriptsize Algo};
        \draw[gde] (p) -- (d); \draw[gde] (d) -- (a); \draw[gde] (m) -- (a);
        \node[font=\scriptsize\bfseries] at (0.95,-2.2) {prerequisites: a DAG};
    \end{scope}
    % maze
    \begin{scope}[xshift=8.6cm, yshift=-1.5cm]
        \foreach \r/\c/\t in {0/0/gcell, 0/1/gcell, 0/2/gwall, 1/0/gwall, 1/1/gcell, 1/2/gcell, 2/0/gcell, 2/1/gcell, 2/2/gwall}
            \node[\t, minimum size=7mm] at (\c*0.7, 2.8-\r*0.7) {};
        \draw[primaryblue, line width=1.2pt] (0,2.8) -- (0.7,2.8) -- (0.7,2.1) -- (1.4,2.1);
        \draw[primaryblue, line width=1.2pt] (0.7,2.1) -- (0.7,1.4) -- (0,1.4);
        \foreach \x/\y in {0/2.8, 0.7/2.8, 0.7/2.1, 1.4/2.1, 0.7/1.4, 0/1.4}
            \fill[primaryblue] (\x,\y) circle (2.2pt);
        \node[font=\scriptsize\bfseries] at (0.7,0.8) {maze: cells + doors};
    \end{scope}
\end{tikzpicture}
\figcap{Three faces of the same object. Friendships (undirected; the two colours show the two separate groups), prerequisites (directed, no cycle), and a maze whose free cells become vertices joined when they share a side (dark cells are walls).}
\end{diagrambox}

%------------------------------------------------------------
\subsubsection*{How a graph arrives in contest input}

The standard format is ``$n$ $m$'' on the first line, followed by $m$ lines with one edge each (and a weight, if the graph is weighted). LeetCode passes the same information as \texttt{n} and a vector \texttt{edges} of pairs. Learn to \emph{draw} such input on paper in under a minute; it is the first thing to do with every sample.

\begin{diagrambox}
\begin{tikzpicture}
    \node[draw=darkgray, fill=lightgray!60, rounded corners=2pt, font=\ttfamily\small, align=left, inner sep=5pt] (in) at (0,0)
        {6 5\\1 2\\2 3\\1 3\\4 5\\5 6};
    \node[font=\scriptsize\color{darkgray}, above=1mm of in] {input};
    \draw[flow, line width=1.2pt] (0.9,0) -- (2.1,0);
    \begin{scope}[xshift=3.0cm, yshift=-0.2cm]
        \node[gv] (1) at (0,0.9) {1};
        \node[gv] (2) at (1.3,0.9) {2};
        \node[gv] (3) at (0.65,-0.3) {3};
        \node[gv] (4) at (2.7,0.9) {4};
        \node[gv] (5) at (4.0,0.9) {5};
        \node[gv] (6) at (3.35,-0.3) {6};
        \foreach \a/\b in {1/2, 2/3, 1/3, 4/5, 5/6} \draw[ge] (\a) -- (\b);
    \end{scope}
    \node[note, text width=4cm] at (10.0,0.3) {$n=6$ vertices, $m=5$ edges. The graph falls into two pieces: $\{1,2,3\}$ and $\{4,5,6\}$.};
\end{tikzpicture}
\figcap{Read the first line, draw $n$ dots, then one line per edge. Nothing says the graph is in one piece --- here it is not.}
\end{diagrambox}

%------------------------------------------------------------
\subsubsection*{Drawing conventions used in this book}

Every diagram in the book uses the same visual vocabulary, so a picture can be read without its caption once you know it. Algorithms colour vertices as they progress; the colours always mean the same thing.

\begin{diagrambox}
\begin{tikzpicture}[x=1cm, y=1cm]
    \foreach \s/\t/\d [count=\i from 0] in {gv/plain/a vertex, gvsrc/source/start here, gvdst/target/answer, gvcur/current/being processed,
                                         gvseen/discovered/queued or stacked, gvdone/finished/fully explored, gvdead/removed/ignored, gvsel/selected/chosen set}{
        \pgfmathsetmacro{\xx}{mod(\i,4)*3.6}
        \pgfmathsetmacro{\yy}{-floor(\i/4)*1.3}
        \node[\s] at (\xx,\yy) {v};
        \node[anchor=west, font=\scriptsize\bfseries] at (\xx+0.5,\yy+0.15) {\t};
        \node[anchor=west, font=\scriptsize\color{darkgray}] at (\xx+0.5,\yy-0.2) {\d};
    }
    \node[gvA] at (0,-2.6) {v}; \node[anchor=west, font=\scriptsize\bfseries] at (0.5,-2.45) {side A};
    \node[anchor=west, font=\scriptsize\color{darkgray}] at (0.5,-2.8) {colour class 1};
    \node[gvB] at (3.6,-2.6) {v}; \node[anchor=west, font=\scriptsize\bfseries] at (4.1,-2.45) {side B};
    \node[anchor=west, font=\scriptsize\color{darkgray}] at (4.1,-2.8) {colour class 2};
    \node[gcell] at (7.2,-2.6) {}; \node[anchor=west, font=\scriptsize\bfseries] at (7.7,-2.45) {grid cell};
    \node[anchor=west, font=\scriptsize\color{darkgray}] at (7.7,-2.8) {free square};
    \node[gwall] at (10.8,-2.6) {}; \node[anchor=west, font=\scriptsize\bfseries] at (11.3,-2.45) {wall};
    \node[anchor=west, font=\scriptsize\color{darkgray}] at (11.3,-2.8) {blocked};
\end{tikzpicture}
\figcap{Vertex and cell styles. A traversal turns vertices from plain to \emph{discovered} (orange) to \emph{current} (yellow) to \emph{finished} (blue).}
\end{diagrambox}

\begin{diagrambox}
\begin{tikzpicture}[x=1cm, y=1cm]
    \foreach \s/\t/\d [count=\i from 0] in {ge/edge/undirected, gde/arc/directed, gtree/tree edge/used by a search, gpath/path/chosen route,
                                         gback/back edge/closes a cycle, gfaded/faded/not relevant now, gcut/cut/separates parts}{
        \pgfmathsetmacro{\xx}{mod(\i,4)*3.6}
        \pgfmathsetmacro{\yy}{-floor(\i/4)*1.1}
        \draw[\s] (\xx-0.4,\yy) -- (\xx+0.5,\yy);
        \node[anchor=west, font=\scriptsize\bfseries] at (\xx+0.6,\yy+0.15) {\t};
        \node[anchor=west, font=\scriptsize\color{darkgray}] at (\xx+0.6,\yy-0.2) {\d};
    }
    \node[gw] at (10.8,-1.1) {5};
    \node[anchor=west, font=\scriptsize\bfseries] at (11.1,-0.95) {weight};
    \node[anchor=west, font=\scriptsize\color{darkgray}] at (11.1,-1.3) {$w(u,v)=5$};
\end{tikzpicture}
\figcap{Edge styles and labels. A small green number next to a vertex, like \tikz[baseline=-0.6ex]\node[gdist]{3};, is a distance computed by an algorithm.}
\end{diagrambox}

%------------------------------------------------------------
\subsubsection*{Check your understanding}

\begin{enumerate}[leftmargin=*]
    \item Draw the graph given by \texttt{n = 4, edges = [[1,2],[2,3],[3,1]]}. How many vertices have no neighbour?
    \item ``Each of the $n$ employees except the director has exactly one direct boss.'' What are the vertices and edges, and should the edges be directed?
    \item Input \texttt{3 4 / 1 2 / 2 1 / 2 3 / 3 3} describes an undirected graph. Is it simple? If not, what makes it a multigraph?
    \item A map has 6 towns and every pair of towns is joined by a two-way road. How many edges? How many arcs if every road is replaced by two one-way roads?
    \item In a weighted graph, is the path with the fewest edges always the cheapest path? Give a reason or a counter-example from this lesson.
\end{enumerate}

\begin{solutionbox}[Answers]
\begin{enumerate}[leftmargin=*]
    \item A triangle 1--2--3 plus the isolated vertex 4: \textbf{one} vertex (4) has no neighbour. Isolated vertices are easy to forget when $n$ is given separately from the edges.
    \item Vertices = employees, one edge per employee--boss pair, $n-1$ edges. Direct them (boss $\to$ employee, or employee $\to$ boss) because ``is the boss of'' is not symmetric. This graph is a \emph{rooted tree} (Lesson~\ref{ssec:ch1-vocabulary}).
    \item Not simple: \texttt{1 2} and \texttt{2 1} are parallel edges and \texttt{3 3} is a self-loop. The program above prints \texttt{self-loops: 1}, \texttt{parallel edges: 1}.
    \item $\binom{6}{2}=15$ undirected edges; $6\cdot 5=30$ arcs.
    \item No. In the weighted figure above, S--A--T and S--B--T have 2 edges and weights 9 and 6, while S--A--B--T has 3 edges and weight 4.
\end{enumerate}
\end{solutionbox}
